\begin{table}[htbp]
  \centering
  \caption{Logical channels and normalized scheduling costs.}
  \label{tab:sensor_specs}
  \small
  \setlength{\tabcolsep}{3.5pt}
  \begin{tabular}{@{}c >{\raggedright\arraybackslash}p{0.27\linewidth} >{\raggedright\arraybackslash}p{0.33\linewidth} c c c@{}}
    \toprule
    Idx & Logical channel & Main observed variables
        & $p_i$ & $p_i^{\mathrm{peak}}$ & Warm-up \\
    \midrule
    1 & Meteorological core channel
      & wind, air temperature, relative humidity, pressure
      & 0.25 & 0.30 & 0 \\
    2 & Basic radiometer channel
      & solar radiation
      & 0.08 & 0.10 & 0 \\
    3 & Shielded temperature--humidity channel
      & air temperature, relative humidity
      & 0.18 & 0.22 & 0 \\
    4 & Infrared snow-surface-temperature channel
      & snow surface temperature
      & 0.49 & 0.62 & 1 \\
    5 & Laser disdrometer channel
      & particle diameter, particle velocity, simulated particle latent variable
      & 0.49 & 0.62 & 2 \\
    6 & FC4 snow mass-flux channel
      & snow mass flux
      & 0.49 & 0.62 & 1 \\
    \bottomrule
  \end{tabular}
  \vspace{0.5ex}

  \noindent\footnotesize
  The meteorological core channel is mandatory. At each time step, the scheduler may
  activate at most one optional channel and may leave all optional channels
  inactive. The infrared snow-surface-temperature channel, laser disdrometer channel,
  and FC4 snow mass-flux channel are
  event-specific channels for thermal, particle, and flux event types,
  respectively; the basic radiometer channel and shielded temperature--humidity channel
  are additional optional channels. Condition labels are simulator-derived metadata
  used only for training, calibration/validation, or offline evaluation; they are not
  direct channel measurements.
\end{table}
